\section{Gauge equivalence for direct sums}%
\label{sec:multi_block_ft}
%% ---------------------------------------------------------

The first lemmas in this section are elementary facts about block-diagonal tensors:
they assemble already-known block gauges, and they record the easy consequences
of applying the injective single-block theorem to each block separately. They do
not derive the block matching appearing in
\cite[Theorem~2.10]{Cirac2016MPDO_arXiv}. In the source proof, that matching
is obtained from BNT overlap, independence, and coefficient comparison; the
corresponding global gauge assembly is
Theorem~\ref{thm:sector_bnt_equal_global_gauge}.

\begin{lemma}[Gauge equivalence of direct sums from per-block gauge equivalence]%
    \label{lem:gauge_equiv_from_block_gauge}
    \lean{MPSTensor.gaugeEquiv_toTensorFromBlocks_of_blockGauge}\leanok
    \uses{lem:tensor_from_blocks_globalGauge_conj, def:gauge_equiv}
    If each pair $(A_k, B_k)$ is gauge equivalent, then the weighted
    block-diagonal tensors $\bigoplus_k \mu_k A_k$ and
    $\bigoplus_k \mu_k B_k$ are gauge equivalent.
\end{lemma}

\begin{proof}\leanok
    Choose per-block gauge matrices via the hypothesis; then apply
    Lemma~\ref{lem:tensor_from_blocks_globalGauge_conj}.
\end{proof}

\begin{lemma}[Gauge equivalence of direct sums with phase-weight absorption]%
    \label{lem:gauge_equiv_block_gauge_phase_weight}
    \lean{MPSTensor.gaugeEquiv_toTensorFromBlocks_of_blockGaugePhase_weight}\leanok
    \uses{lem:tensor_from_blocks_globalGauge_conj, def:gauge_equiv,
        def:gauge_phase_equiv}
    Let $A_k, B_k$ be block families, and let $X_k \in \GL_{D_k}(\C)$
    and scalars $\zeta_k \in \C$ satisfy
    $B_k^i = \zeta_k \cdot X_k A_k^i X_k^{-1}$ for every block $k$ and
    physical index $i$. If the weights satisfy $\mu_A(k) = \mu_B(k) \zeta_k$,
    then the weighted block-diagonal tensors
    $\bigoplus_k \mu_A(k) A_k$ and $\bigoplus_k \mu_B(k) B_k$ are gauge
    equivalent.
\end{lemma}

\begin{proof}\leanok
    \uses{lem:tensor_from_blocks_globalGauge_conj}
    Per-block, the weight identity $\mu_A(k) = \mu_B(k) \zeta_k$ rewrites the
    scaled relation as
    $\mu_B(k)\,B_k^i = X_k\,(\mu_A(k)\,A_k^i)\,X_k^{-1}$.
    Each weighted block then satisfies an ordinary conjugation, and the
    block-diagonal conjugation formula assembles the per-block $X_k$ into the
    flattened global gauge $\bigoplus_k X_k$ relating
    $\bigoplus_k \mu_A(k) A_k$ to $\bigoplus_k \mu_B(k) B_k$.
\end{proof}

\subsection{Same-index direct-sum corollaries from the single-block theorem}%
\label{subsec:ft_block_sum_corollaries}

\begin{remark}
    See Definition~\ref{def:block_diagonal_tensor} for the
    block-diagonal tensor construction. The results in this subsection assume
    that the block correspondence has already been fixed; they do not prove
    \cite[Theorem~2.10]{Cirac2016MPDO_arXiv}. They are same-index
    direct-sum corollaries, not BNT preparation steps in the source proof
    or in the after-blocking canonical-form reduction.
\end{remark}

\begin{lemma}[Per-block application of the single-block theorem]\label{lem:multi_block_ft}
    \lean{MPSTensor.fundamentalTheorem_multiBlock_blocks}
    \leanok
    \uses{def:block_diagonal_tensor, def:injective}
    If all block tensors $A_k$ are injective and the per-block
    same-MPV condition holds ($A_k$ and $B_k$ generate the same
    MPV family for each~$k$),
    then each pair $(A_k, B_k)$ is gauge equivalent.
\end{lemma}

\begin{proof}
    \leanok
    \uses{thm:ft_single}
    Apply the single-block Fundamental Theorem
    (Theorem~\ref{thm:ft_single}) to each block independently.
\end{proof}

\begin{remark}[Per-block restatement, not the multi-block paper theorem]
    This is a block-sum reformulation of the per-block injective single-block FT:
    the block pairing $(A_k, B_k)$ is assumed already given and the blocks already
    matched. It is not the multi-block Fundamental Theorem of
    \cite[Theorem~2.10, lines~349--352]{Cirac2016MPDO_arXiv}, which derives the
    block matching and permutation from only the global same-MPV hypothesis.
\end{remark}

\begin{lemma}[Per-block same MPV implies global same MPV]\label{lem:block_to_global_mpv}
    \lean{MPSTensor.sameMPV_toTensorFromBlocks_of_blockSameMPV}
    \leanok
    \uses{def:block_diagonal_tensor, def:same_mpv}
    If $A_k$ and $B_k$ generate the same MPV family for each
    block~$k$, then the block-diagonal tensors
    $\bigoplus_k \mu_k A_k$ and $\bigoplus_k \mu_k B_k$
    also generate the same MPV family.
\end{lemma}

\begin{proof}\leanok\uses{thm:mpv_decomposition}
    By the MPV decomposition (Theorem~\ref{thm:mpv_decomposition}),
    each MPV is $\sum_k \mu_k^N V^{(N)}(A_k)_\sigma$.
    Per-block equality of MPV coefficients gives global equality.
\end{proof}

\begin{remark}[Direct-sum congruence]
    This is a direct-sum congruence statement: per-block MPV equality lifts to
    global MPV equality via the block-diagonal decomposition formula.
\end{remark}

\begin{lemma}[Gauge equivalence for same-index direct sums]\label{lem:multi_block_global}
    \lean{MPSTensor.fundamentalTheorem_multiBlock_global}
    \leanok
    \uses{def:block_diagonal_tensor, def:injective, def:same_mpv, lem:multi_block_ft}
    If each block $A_k$ is injective and each pair $(A_k, B_k)$ generates
    the same MPV family, then the block-diagonal tensors
    $\bigoplus_k \mu_k A_k$ and $\bigoplus_k \mu_k B_k$ are gauge equivalent.
\end{lemma}

\begin{proof}\leanok\uses{lem:multi_block_ft}
    Apply the per-block single-block Fundamental Theorem
    (Lemma~\ref{lem:multi_block_ft}) to obtain blockwise gauge
    equivalences from injectivity and equality of MPV families, then
    combine them into a global gauge for the block-diagonal tensors.
\end{proof}

\begin{remark}[Block-diagonal gauge assembly, not the FT block-matching theorem]
    This combines per-block gauges into a block-diagonal gauge for the global
    assembled tensor, given pre-matched same-index blocks. It is not the FT
    block-matching theorem: it does not derive the block pairing or permutation
    from a global same-MPV hypothesis; the block correspondence is assumed in
    advance.
\end{remark}

\begin{lemma}[Gauge equivalence of direct sums from per-block MPV equality]%
    \label{lem:fundamentalTheorem_multiBlock_full}
    \lean{MPSTensor.fundamentalTheorem_multiBlock_full}
    \leanok
    \uses{def:injective, def:same_mpv, def:gauge_equiv,
        def:block_diagonal_tensor, lem:tensor_from_blocks_globalGauge_conj}
    Let $\{A_k\}_{k=0}^{r-1}$ and $\{B_k\}_{k=0}^{r-1}$ be two families
    of MPS tensors with block dimensions $D_k$, each $A_k$ injective, and
    $\mc{V}(A_k) = \mc{V}(B_k)$ for all~$k$.
    Then:
    \begin{enumerate}
        \item Per-block gauge equivalence:
              $B_k^i = X_k\, A_k^i\, X_k^{-1}$ for some
              $X_k \in \GL_{D_k}(\C)$, for every $k$ and $i$.
        \item Global gauge equivalence: the block-diagonal tensors
              $\bigoplus_k \mu_k A_k$ and $\bigoplus_k \mu_k B_k$ are
              gauge equivalent.
    \end{enumerate}
    Explicitly, for $X=\bigoplus_k X_k$,
    \begin{align}
        \bigoplus_k \mu_k B_k^i
         & = X \left(\bigoplus_k \mu_k A_k^i\right) X^{-1}.
        \label{eq:aft_multiBlock_global_conj}
    \end{align}
\end{lemma}

\begin{proof}\leanok\uses{lem:multi_block_ft, lem:multi_block_global}
    Per-block gauge equivalence follows from applying the single-block
    Fundamental Theorem to each block, using injectivity of
    $A_k$ and $\mc{V}(A_k) = \mc{V}(B_k)$.
    With $X=\bigoplus_k X_k$, (\ref{eq:aft_multiBlock_global_conj})
    follows by multiplying block-diagonal matrices block by block.
\end{proof}

\begin{lemma}[Explicit per-block gauges]%
    \label{lem:fundamentalTheorem_multiBlock_explicit}
    \lean{MPSTensor.fundamentalTheorem_multiBlock_explicit}
    \leanok
    \uses{lem:fundamentalTheorem_multiBlock_full}
    Under the hypotheses of
    Lemma~\ref{lem:fundamentalTheorem_multiBlock_full}, there exist
    invertible matrices $X_k \in \GL_{D_k}(\C)$ such that
    $B_k^i = X_k\, A_k^i\, X_k^{-1}$ for every block $k$ and every
    physical index $i$.
\end{lemma}

\begin{proof}
    \leanok
    \uses{lem:fundamentalTheorem_multiBlock_full}
    Extract the gauge matrix from the per-block part of
    Lemma~\ref{lem:fundamentalTheorem_multiBlock_full} for each block.
\end{proof}

\begin{lemma}[Product-algebra decomposition of per-block MPV equality]%
    \label{lem:fundamentalTheorem_multiBlock_decomposition}
    \lean{MPSTensor.fundamentalTheorem_multiBlock_decomposition}
    \leanok
    \uses{lem:fundamentalTheorem_multiBlock_full}
    Assume that every block dimension is nonzero.
    Under the hypotheses of
    Lemma~\ref{lem:fundamentalTheorem_multiBlock_full}, the product-algebra
    automorphism determined by the per-block linear extensions decomposes into
    a block permutation together with inner conjugation on each block: there
    exist a permutation $\sigma$ of the blocks, compatible block-dimension
    equalities, and invertible matrices $X_k$ such that, after reindexing the
    target block $\sigma(k)$ back to block $k$, the action on the $k$th matrix
    block is $M \mapsto X_k M X_k^{-1}$.
\end{lemma}

\begin{proof}
    \leanok
    \uses{thm:piAlgEquiv_decomposition}
    Apply Theorem~\ref{thm:piAlgEquiv_decomposition} to the product-algebra
    automorphism built from the per-block linear extensions.
\end{proof}

\subsection{MPV invariance under reindexing}%
\label{subsec:mpv_reindexing}

\begin{lemma}[MPV invariance under block permutation]\label{lem:mpv_perm_blocks}
    \lean{MPSTensor.sameMPV₂_toTensorFromBlocks_perm}\leanok
    \uses{def:block_diagonal_tensor, def:same_mpv2}
    Permuting the block index set and transporting the weights
    correspondingly does not change the $\mathrm{SameMPV}_2$ family of the
    weighted block-diagonal tensor.
\end{lemma}

\begin{proof}\leanok
    \uses{def:block_diagonal_tensor}
    The sum of per-block MPV contributions is invariant under rearrangement
    of the summands.
\end{proof}

\begin{lemma}[MPV invariance under dimension cast]\label{lem:mpv_cast_blocks}
    \lean{MPSTensor.sameMPV₂_toTensorFromBlocks_cast}\leanok
    \uses{def:block_diagonal_tensor, def:same_mpv2}
    If two block-dimension families agree pointwise, the corresponding
    weighted block-diagonal tensors generate the same $\mathrm{SameMPV}_2$
    family after transporting the blocks across the dimension equalities.
\end{lemma}

\begin{proof}\leanok
    Substituting the equality of dimensions makes the two tensor families
    identical.
\end{proof}

\begin{lemma}[MPV invariance under block reindexing]\label{lem:mpv_reindex_blocks}
    \lean{MPSTensor.mpv_toTensorFromBlocks_reindex}\leanok
    \uses{def:block_diagonal_tensor, thm:mpv_decomposition}
    Let $e : \{0,\ldots,r'-1\}\simeq\{0,\ldots,r-1\}$ be a bijection of block
    indices. For weights $\mu_k$, block tensors $A_k$, word length $N$, and
    configuration $\sigma$, the coefficient of the block-diagonal MPV is unchanged
    by reindexing the blocks along $e$:
    \begin{align}
        V^{(N)}\!\left(\bigoplus_k \mu_k A_k\right)_\sigma
         & =
        V^{(N)}\!\left(\bigoplus_k \mu_{e(k)} A_{e(k)}\right)_\sigma.
        \notag
    \end{align}
\end{lemma}

\begin{proof}\leanok
    \uses{thm:mpv_decomposition}
    Expanding both sides by Theorem~\ref{thm:mpv_decomposition} gives
    $\sum_k \mu_k^N\,V^{(N)}(A_k)_\sigma$ and
    $\sum_k \mu_{e(k)}^N\,V^{(N)}(A_{e(k)})_\sigma$.
    These sums are equal by reindexing the finite sum along the bijection $e$.
\end{proof}

\subsection{Converse: gauge equivalence implies same MPV}%
\label{subsec:converse_gauge_mpv}

\begin{lemma}[Gauge equivalence of direct sums implies same MPV]\label{lem:gauge_equiv_implies_same_mpv}
    \lean{MPSTensor.gaugeEquiv_toTensorFromBlocks_implies_sameMPV}\leanok
    \uses{def:block_diagonal_tensor, def:gauge_equiv, def:same_mpv}
    If the weighted block-diagonal tensors $\bigoplus_k \mu_k A_k$ and
    $\bigoplus_k \mu_k B_k$ are gauge equivalent, then they generate the
    same MPV family.
\end{lemma}

\begin{proof}\leanok
    \uses{def:gauge_equiv}
    Gauge equivalence preserves all traces of word evaluations, hence all
    MPV coefficients.
\end{proof}

\subsection{Canonical-form tensor as a weighted direct sum}%
\label{subsec:canonical_form_block_sum}

\begin{lemma}[Canonical-form tensor agrees with block sum]\label{lem:canonical_toTensor_eq}
    \lean{MPSTensor.CanonicalForm.toTensor_eq_toTensorFromBlocks}\leanok
    \uses{def:block_diagonal_tensor}
    For a canonical-form record $C$, the assembled tensor $T_C$ is
    definitionally equal to the weighted block-diagonal tensor
    $\bigoplus_k \mu_k^C\, A_k^C$,
    where $\mu_k^C$ are the block weights and $A_k^C$ are the block tensors
    stored in $C$.
\end{lemma}

\begin{proof}\leanok
    The equality holds by definition; $T_C$ is defined as the weighted
    block-diagonal tensor formed from the weights and block tensors of $C$.
\end{proof}

\begin{theorem}[Fundamental theorem for canonical forms with the same block structure]\label{thm:ft_canonical_form_same_structure}
    \lean{MPSTensor.fundamentalTheorem_canonicalForm_sameStructure}\leanok
    \uses{lem:canonical_toTensor_eq, lem:multi_block_global}
    Let $C$ be a canonical-form record with block weights $\mu_k^C$ and block
    tensors $A_k^C$. Let $\{B_k\}$ be a family of MPS tensors on the same
    block dimensions, and assume that each $A_k^C$ is injective and generates
    the same MPV family as $B_k$.
    Then $T_C$ and $\bigoplus_k \mu_k^C\, B_k$ are gauge equivalent.
\end{theorem}

\begin{proof}\leanok
    \uses{lem:canonical_toTensor_eq, lem:multi_block_global}
    Rewrite $T_C$ as a weighted block-diagonal tensor via
    Lemma~\ref{lem:canonical_toTensor_eq}, then apply the multi-block
    global gauge-equivalence lemma
    (Lemma~\ref{lem:multi_block_global}).
\end{proof}

\begin{lemma}[Per-block equality of MPV families and per-block gauge equivalence]%
    \label{lem:perBlock_sameMPV_iff_gaugeEquiv}
    \lean{MPSTensor.perBlock_sameMPV_iff_gaugeEquiv}
    \leanok
    \uses{def:injective, def:same_mpv, def:gauge_equiv}
    Let $\{A_k\}$ be a family of injective block tensors.
    Then $\mc{V}(A_k) = \mc{V}(B_k)$ for all~$k$ if and only if
    $A_k$ and $B_k$ are gauge equivalent for all~$k$.
\end{lemma}

\begin{proof}\leanok\uses{def:gauge_equiv}
    The forward direction applies the single-block Fundamental Theorem to
    each block.
    The reverse direction holds because gauge equivalence preserves all
    traces of word evaluations.
\end{proof}

\subsection{Equal-case block matching with differing bond dimensions}%
\label{sec:different_bond_dimension_equal_case_ft}

The same-index lemmas above assume a fixed block correspondence. The equal-case
sector matching with initially different sector counts and bond dimensions is
Theorem~\ref{lem:ft_sector_bnt_equal_sector_data_different_dims} of
Chapter~\ref{ch:ft_proof}; its global-gauge consequence is
Theorem~\ref{thm:sector_bnt_equal_global_gauge}.

%% ---------------------------------------------------------
\section{Translation-invariant reduction and global symmetry}%
\label{sec:ti_reduction}
%% ---------------------------------------------------------

The single-gauge corollary for translation-invariant chains specializes to
a reduction statement: equal states force one tensor to be a gauge of the
other and inherit injectivity. Composing two gauges of the same tensor then
yields the scalar law underlying the global-symmetry consequences.

\begin{theorem}[Translation-invariant reduction corollary]%
    \label{thm:ti_reduction_corollary}
    \lean{MPSChainTensor.ti_reduction_corollary}
    \leanok
    \uses{def:injective, def:same_mpv}
    Let $A$ and $B$ be MPS tensors of bond dimension $D$ with $A$
    injective.
    Consider the two constant translation-invariant chains
    $(A, \ldots, A)$ and $(B, \ldots, B)$ of length $n \ge 1$.
    If the combined tensors of the two chains generate the same MPV
    family, then:
    \begin{enumerate}
        \item $B$ is gauge equivalent to $A$: there exists
              $X \in \GL_D(\C)$ with $B^i = X\, A^i\, X^{-1}$ for
              all~$i$.
        \item $B$ is injective.
    \end{enumerate}
\end{theorem}

\begin{proof}\leanok\uses{thm:ti_tensors_single_gauge}
    Theorem~\ref{thm:ti_tensors_single_gauge} gives a single gauge matrix
    $X$ with $B^i = X A^i X^{-1}$ for all $i$.
    Since gauge equivalence preserves injectivity, $B$ is injective as well.
\end{proof}

\subsection{Global symmetry via gauge composition}

The virtual representation theorem
(Theorem~\ref{thm:virtual_rep_injective}) relies on the fact that
composing two gauge transformations yields another gauge up to a
scalar factor.

\begin{lemma}[Gauge ratio commutes with tensor matrices]%
    \label{lem:gauge_ratio_commutes}
    \lean{MPSTensor.gauge_ratio_commutes}
    \leanok
    \uses{def:gauge_equiv}
    Diagrammatically, the hypothesis is that the two local gauges
    \begin{center}
        % B^i = X A^i X^{-1}: the red capsules act on the west and east
        % virtual indices of A^i; the upper leg is the free index i.
        % Both sides have boundary signature (1,1,1,0).
        \begin{tenkz}[physical=up]
            \tn[up=$i$]{B}
        \end{tenkz}
        $ = $
        \begin{tenkz}[physical=up]
            \tnX{X} & \tn[up=$i$]{A} & \tnX{X^{-1}}
        \end{tenkz}

        \qquad\text{and}\qquad

        % B^i = Y A^i Y^{-1}: this second gauge has the same source,
        % target, and free physical index as the X-gauge above.  Both
        % sides again have boundary signature (1,1,1,0).
        \begin{tenkz}[physical=up]
            \tn[up=$i$]{B}
        \end{tenkz}
        $ = $
        \begin{tenkz}[physical=up]
            \tnX{Y} & \tn[up=$i$]{A} & \tnX{Y^{-1}}
        \end{tenkz}
    \end{center}
    produce the same target tensor $B^i$ from the same source tensor
    $A^i$.
    If $X, Y \in \GL_D(\C)$ both conjugate an MPS tensor $A$ to the same
    tensor $B$ (i.e.\ $B^i = X A^i X^{-1} = Y A^i Y^{-1}$ for all $i$),
    then $Y^{-1} X$ commutes with every $A^i$.
\end{lemma}

\begin{proof}\leanok
    Direct matrix calculation: conjugating both gauge relations and
    composing gives $(Y^{-1} X) A^i = A^i (Y^{-1} X)$.
\end{proof}

\begin{lemma}[Gauge ratio is scalar for injective tensors]%
    \label{lem:gauge_ratio_isScalar}
    \lean{MPSTensor.gauge_ratio_isScalar}
    \leanok
    \uses{lem:gauge_ratio_commutes, def:injective, lem:center_scalar}
    Under the same hypotheses, if $A$ is additionally injective, then
    $Y^{-1} X = \lambda \Id_D$ for some scalar $\lambda \in \C$.
\end{lemma}

\begin{proof}\leanok\uses{lem:gauge_ratio_commutes, lem:center_scalar}
    The same two-gauge picture shows that the ratio $Y^{-1}X$ acts
    invisibly on the local tensor.
    By Lemma~\ref{lem:gauge_ratio_commutes}, $Y^{-1}X$ commutes with
    all $A^i$.
    Since $\{A^i\}$ spans $\MN{D}$, the matrix $Y^{-1}X$ commutes
    with all of $\MN{D}$.
    By Lemma~\ref{lem:center_scalar}, it is a scalar matrix.
\end{proof}

\begin{theorem}[Projective multiplication law for gauge matrices]%
    \label{thm:gaugeMatrix_projective_mul}
    \lean{MPSTensor.gaugeMatrix_projective_mul}
    \leanok
    \uses{lem:gauge_ratio_isScalar, def:injective,
        def:on_site_symmetric, lem:twisted_tensor_mul}
    Let $A$ be an injective MPS tensor with on-site symmetry under a
    group representation $U : G \to \GL_d(\C)$.
    For each $g \in G$, let $X(g) \in \GL_D(\C)$ satisfy
    $\widetilde{A}_g^i = X(g)\, A^i\, X(g)^{-1}$ for all~$i$.
    Then for all $g, h \in G$, there exists $c(g,h) \in \C$ such that
    $X(h)\, X(g) = c(g,h)\, X(g \cdot h)$.
    In fact $c(g,h) \ne 0$ since both sides are invertible, so
    $c(g,h) \in \C^{\times}$. Diagrammatically, the two candidate gauges
    are competing local replacements for the same twisted tensor:
    \begin{center}
        % (X(h)X(g)) A^i (X(h)X(g))^{-1}
        %   = X(gh) A^i X(gh)^{-1}: the two candidate gauges act on
        % the same virtual indices of A^i and produce the same twisted
        % tensor.  The upper leg is i; both terms have signature (1,1,1,0).
        \begin{tenkz}[physical=up]
            \tnX{X(h)X(g)} & \tn[up=$i$]{A} &
            \tnX{(X(h)X(g))^{-1}}
        \end{tenkz}
        $ = $
        \begin{tenkz}[physical=up]
            \tnX{X(gh)} & \tn[up=$i$]{A} & \tnX{X(gh)^{-1}}
        \end{tenkz}
    \end{center}
    and both sides conjugate $A^i$ to the same twisted tensor
    $\widetilde{A}_{gh}^i$.
\end{theorem}

\begin{proof}\leanok
    \uses{lem:gauge_ratio_isScalar, lem:twisted_tensor_mul}
    By the composition law $\widetilde{A}_{gh} = \widetilde{(\widetilde{A}_h)}_g$
    (Lemma~\ref{lem:twisted_tensor_mul}), both $X(h)\,X(g)$ and $X(gh)$
    conjugate $A$ to $\widetilde{A}_{gh}$.
    Lemma~\ref{lem:gauge_ratio_isScalar} then gives the scalar factor.
\end{proof}
